\section*{Original Section 3.3: Explicit Constructions of Variety Evasive Subspace Families (journal pp.32--40)}
Let $n,d\in\mathbb N^+$, $k\in\{0,1,\ldots,n-1\}$, and $\epsilon\in(0,1)$. In this subsection, we prove the Main theorem (Theorem 1.7) by constructing explicit projective or affine $k$-subspace families that are $(n,d,\epsilon)$-evasive.
We first prove Theorem 1.7 in the projective case and then derive the affine case from it by viewing $\Aspace^n$ as an open subset of $\PP^n$. For the projective case, we present two constructions. The first one is simple and only uses $\epsilon$-hitting sets for low degree polynomials (Lemma 2.3). But the size of the resulting subspace family is polynomial only when both $d$ and $k$ (or $n-k$) are bounded. Next, we give a more sophisticated construction, which yields subspace families of polynomial size as long as $d$ is bounded.
\subsection*{3.3.1 Simple Construction}
We first present a simple construction of $(n,d,\epsilon)$-evasive $k$-subspace families on $\PP^n$.
First assume $k+1\le n-k$. In this case, construct a $k$-subspace family $\mathcal H$ on $\PP^n$ as follows:
\begin{enumerate}
\item Use Lemma 2.3 to compute an $\epsilon$-hitting set $T$ for the family of polynomials $f\in\FF[Y_{i,j}:i\in[k+1],j\in[n+1]]$ of degree at most $(k+1)d$ such that $|T|=\operatorname{poly}(\binom{(k+1)(n+1+d)}{(k+1)d},1/\epsilon)$. Think of $T$ as a collection of $(k+1)\times(n+1)$ matrices over $\FF$.
\item Initialize $\mathcal H=\emptyset$. For each matrix $A\in T$, if $A$ has full row rank $k+1$, add to $\mathcal H$ the $k$-subspace $W\in\GG(k,n)$ with the generating matrix $A$.
\end{enumerate}
Next, assume $k+1>n-k$. In this case, construct $\mathcal H$ in a similar way, but use parity check matrices instead of generating matrices. Namely, compute an $\epsilon$-hitting set $T$ for the family of polynomials $f\in\FF[Y_{i,j}:i\in[n-k],j\in[n+1]]$ of degree at most $(n-k)d$ such that $|T|=\operatorname{poly}(\binom{(n-k)(n+1+d)}{(n-k)d},1/\epsilon)$. Think of $T$ as a collection of $(n-k)\times(n+1)$ matrices over $\FF$. For each matrix $A\in T$, add to $\mathcal H$ the $k$-subspace $W\in\GG(k,n)$ with the parity check matrix $A$.
This construction does give an $(n,d,\epsilon)$-evasive $k$-subspace family, as stated by the following lemma.
\begin{result}{Lemma 3.11}
The $k$-subspace family $\mathcal H$ constructed above is $(n,d,\epsilon)$-evasive and has size polynomial in
\[\min\left\{\binom{(k+1)(n+1+d)}{(k+1)d},\binom{(n-k)(n+1+d)}{(n-k)d}\right\}\]
and $1/\epsilon$. Moreover, the total time complexity of computing the linear equations defining the $k$-subspaces in $\mathcal H$ is polynomial in $|\mathcal H|$ (and $\log p$, if $\operatorname{char}(\FF)=p>0$).
\end{result}
\begin{proof}
We only show that $\mathcal H$ is $(n,d,\epsilon)$-evasive since the rest of the lemma is obvious from the construction. Let $\mathcal F$ be the family of all equidimensional subvarieties of $\PP^n$ of dimension $n-k-1$ and degree at most $d$. By Lemma 3.1, it suffices to prove that $\mathcal H$ is $(\mathcal F,\epsilon)$-evasive. Consider any $\mathcal V\in\mathcal F$. We want to show that $\mathcal V\cap W=\emptyset$ for all but at most $\epsilon$-fraction of $W\in\mathcal H$.
First assume $k+1\le n-k$, or equivalently, $\binom{(k+1)(n+1+d)}{(k+1)d}\le\binom{(n-k)(n+1+d)}{(n-k)d}$. The Chow form $\widetilde R_{\mathcal V}$ of $\mathcal V$ in Stiefel coordinates is a nonzero homogeneous polynomial in $\FF[Y_{i,j}:i\in[k+1],j\in[n+1]]$ of degree $(k+1)\deg(\mathcal V)\le(k+1)d$. By the choice of $T$, for all but at most $\epsilon$-fraction of $A\in T$, we have $\widetilde R_{\mathcal V}(A)\ne0$, which implies $\mathcal V\cap W=\emptyset$, where $A$ is a generating matrix of $W$.
By construction, $\mathcal H$ is the collection of $k$-subspaces corresponding to the matrices $A\in T$ of full row rank. So we have ignored the matrices that do not have full row rank. But this does not increase the fraction of ``bad'' $W\in\mathcal H$ since if $A$ does not have full row rank, then the maximal minors of $A$ are all zero, and $\widetilde R_{\mathcal V}(A)$ must be zero. It follows that $\mathcal V\cap W=\emptyset$ for all but at most $\epsilon$-fraction of $W\in\mathcal H$, as desired.
Now assume $k+1>n-k$. The proof in this case is similar and we omit the details. The only difference is that we use the Chow form $\widetilde R^\vee_{\mathcal V}$ in dual Stiefel coordinates instead of $\widetilde R_{\mathcal V}$.
\end{proof}
\subsection*{3.3.2 Improved Construction}
Before presenting the improved construction, we first introduce some notions from algebraic geometry.
\paragraph{Projections.}
Suppose $W$ is a $k$-subspace of $\PP^n$, and $\ell_1,\ldots,\ell_{n-k}\in\FF[X_1,\ldots,X_{n+1}]$ are $n-k$ homogeneous linear polynomials such that $W=V(\ell_1,\ldots,\ell_{n-k})$. Then we have a map $\pi:\PP^n\setminus W\to\PP^{n-k-1}$ defined by
\[\pi:x\mapsto(\ell_1(x),\ldots,\ell_{n-k}(x))\]
which is well-defined since $\ell_1,\ldots,\ell_{n-k}$ never simultaneously vanish on $\PP^n\setminus W$. We say $\pi$ is a projection from $\PP^n\setminus W$ to $\PP^{n-k-1}$ and $W$ is its center. Note that if we lift $\pi$ to the linear map $\pi':\Aspace^{n+1}\to\Aspace^{n-k}$ sending $x\in\Aspace^{n+1}$ to $(\ell_1(x),\ldots,\ell_{n-k}(x))\in\Aspace^{n-k}$, then the center $W$ is simply $\PP\ker(\pi')$.
We need the following lemma.
\begin{result}{Lemma 3.12 (Shafarevich 1994, Section I.5.3, Theorem 7)}
Suppose $\pi:\PP^n\setminus W\to\PP^m$ is a projection with center $W$ and $\mathcal V$ is a subvariety of $\PP^n$ disjoint from $W$. Then $\pi$ restricts to a finite morphism from $\mathcal V$ to $\PP^m$.
\end{result}
\paragraph{Nondegenerate varieties.}
For a subvariety $\mathcal V\subseteq\PP^n$, denote by $\operatorname{span}(\mathcal V)$ the smallest projective subspace that contains $\mathcal V$. We say $\mathcal V$ is nondegenerate if it is not contained in a hyperplane of $\PP^n$, or equivalently, $\operatorname{span}(\mathcal V)=\PP^n$.
We need the following fact from algebraic geometry (see, e.g., Eisenbud \& Harris (1987, Proposition 0) or Harris (1992, Corollary 18.12)).
\begin{result}{Lemma 3.13}
The codimension of a nondegenerate irreducible subvariety $\mathcal V$ of $\PP^n$ is at most $\deg(\mathcal V)-1$.
\end{result}
\paragraph{A two-step construction.}
We now give an improved construction of $(n,d,\epsilon)$-evasive $k$-subspace families on $\PP^n$ as follows.
\begin{enumerate}
\item If $k\le d-2$, just use the previous simple construction. So assume $k>d-2$. Let $k'=d-2<k$, $n'=k'+n-k<n$, and $\epsilon_0=\epsilon/(2d)$.
\item Use Corollary 2.5 to construct a collection $\mathcal H_1$ of $(n'+1)\times(n+1)$ matrix over $\FF$ such that $|\mathcal H_1|=\operatorname{poly}(n,d/\epsilon)$ and for every $(n+1)\times(n'+1)$ matrix $M$ over $\FF$ of rank $n'+1$, all but at most $\epsilon_0$-fraction of $B\in U$ satisfies $\operatorname{rank}(BM)=n'+1$.
We abuse the notation and view $\mathcal H_1$ as a collection of linear maps from $\Aspace^{n+1}$ to $\Aspace^{n'+1}$. Then for any linear subspace $W\subseteq\Aspace^{n+1}$ of dimension $n'+1$, we have $\dim(\pi(W))=\dim(W)=n'+1$ for all but at most $\epsilon_0$-fraction of $\pi\in\mathcal H_1$.
\item Construct a collection $\mathcal H_2$ of linear maps from $\Aspace^{n'+1}$ to $\Aspace^{n-k}$ as follows. First assume $d>1$. Use Lemma 3.11 to construct an $(n',d,\epsilon_0)$-evasive $k'$-subspace family $\mathcal H'_2$ on $\PP^{n'}$ of size polynomial in
\[\min\left\{\binom{(k'+1)(n'+1+d)}{(k'+1)d},\binom{(n-k)(n'+1+d)}{(n-k)d}\right\}\]
and $1/\epsilon_0$. For each $k'$-subspace $W\in\mathcal H_2$, compute a surjective linear map $\pi_W:\Aspace^{n'+1}=\Aspace^{n-k+k'+1}\to\Aspace^{n-k}$ such that $W=\PP\ker(\pi_W)$. Let $\mathcal H_2=\{\pi_W:W\in\mathcal H'_2\}$.
If $d=1$, just let $\mathcal H_2$ be the singleton consisting of the identity map on $\Aspace^{n'+1}=\Aspace^{n-k}$.
\item Initialize $\mathcal H=\emptyset$. For each $(\pi_1,\pi_2)\in\mathcal H_1\times\mathcal H_2$, if $\dim(\ker(\pi_2\circ\pi_1))=k+1$, add the $k$-subspace $\PP\ker(\pi_2\circ\pi_1)$ to $\mathcal H$.
\end{enumerate}
We use the construction above to prove the Main theorem (Theorem 1.7) in the projective case. For convenience, we restate it in the following form.
\begin{result}{Theorem 3.14 (Main theorem in the projective case)}
The $k$-subspace family $\mathcal H$ constructed above is $(n,d,\epsilon)$-evasive and has size $\operatorname{poly}(N(k,d,n),n,1/\epsilon)$. Moreover, the total time complexity of computing the linear equations defining the $k$-subspaces in $\mathcal H$ is polynomial in $|\mathcal H|$ (and $\log p$, if $\operatorname{char}(\FF)=p>0$).
\end{result}
\begin{proof}
The theorem follows from Lemma 3.11 if $k\le d-2$. So assume $k>d-2$. We only show that $\mathcal H$ is $(n,d,\epsilon)$-evasive since the rest of the theorem is obvious from the construction.
Let $\mathcal F$ be the family of all irreducible subvarieties of $\PP^n$ of dimension $n-k-1$ and degree at most $d$. By Lemma 3.3, it suffices to prove that $\mathcal H$ is $(\mathcal F,2\epsilon_0)$-evasive. Consider any $\mathcal V\in\mathcal F$. We want to show that $\mathcal V\cap W=\emptyset$ for all but at most $(2\epsilon_0)$-fraction of $W\in\mathcal H$.
By definition, $\mathcal V$ is a nondegenerate irreducible subvariety of $\operatorname{span}(\mathcal V)$. By Lemma 3.13, the codimension of $\mathcal V$ in $\operatorname{span}(\mathcal V)$ is at most $d-1$. Therefore,
\[\dim(\operatorname{span}(\mathcal V))\le\dim(\mathcal V)+d-1=(n-k-1)+(d-1)=n'.\]
Let $\Lambda\subseteq\PP^n$ be an $n'$-subspace that contains $\operatorname{span}(\mathcal V)$. By the choice of $\mathcal H_1$, all but at most $\epsilon_0$-fraction of $\pi_1\in\mathcal H_1$ satisfies $\PP\ker(\pi_1)\cap\Lambda=\emptyset$. Fix $\pi_1$ that satisfies this condition. Then $\PP\ker(\pi_1)$ is disjoint from $\mathcal V\subseteq\Lambda$. By Lemma 3.12, $\pi_1$ induces a finite morphism $\overline\pi_1:\mathcal V\to\PP^{n'}$.
Let $\mathcal V'=\overline\pi_1(\mathcal V)\subseteq\PP^{n'}$. Then $\mathcal V'$ is a projective subvariety of dimension $\dim(\mathcal V)=n-k-1$ and degree at most $d$. (The degree bound follows from, e.g., an inductive application of Mumford (1976, Proposition 5.5).) By the choice of $\mathcal H_2$, all but at most $\epsilon_0$-fraction of $\pi_2\in\mathcal H_2$ satisfies $\PP\ker(\pi_2)\cap\mathcal V'=\emptyset$. Fix $\pi_2$ that satisfies this condition. (If $d=1$ and $\pi_2$ is the identity map, we regard $\PP\ker(\pi_2)$ as an empty set, in which case this condition is also satisfied.) By Lemma 3.12, $\pi_2$ induces a finite morphism $\overline\pi_2:\mathcal V'\to\PP^{n-k-1}$. So we have a finite morphism $\overline\pi_2\circ\overline\pi_1:\mathcal V\to\PP^{n-k-1}$. Note that $\overline\pi_2\circ\overline\pi_1$ is defined by restricting a projection with center $\PP\ker(\pi_2\circ\pi_1)$ to $\mathcal V$. As $\overline\pi_2\circ\overline\pi_1$ is well-defined on $\mathcal V$, its center $\PP\ker(\pi_2\circ\pi_1)$ is disjoint from $\mathcal V$ and this also forces $\dim\ker(\pi_2\circ\pi_1)=k+1$ by Lemma 2.9.
By the above argument and the construction of $\mathcal H$, all but at most $2\epsilon_0$-fraction of the $k$-subspaces in $\mathcal H$ are disjoint from $\mathcal V$, as desired.
\end{proof}
\section*{Original Section 3.3.4: Proof of the Main theorem in the Affine Case (journal pp.39--40)}
We now prove Theorem 1.7 in the affine case. Recall that we may view $\Aspace^n$ as an open subset of $\PP^n$ via the map $(x_1,\ldots,x_n)\mapsto(x_1,\ldots,x_n,1)$. In this way, $\PP^n$ becomes the disjoint union of $\Aspace^n$ and the hyperplane at infinity $H_\infty$ defined by $X_{n+1}=0$.
We use the following lemma to reduce the affine case to the projective case.
\begin{result}{Lemma 3.15}
Let $n,d\in\mathbb N^+$, $k\in\{0,1,\ldots,n-1\}$, and $\epsilon'\in(0,1/2)$. Suppose $\mathcal H$ is an $(n,d,\epsilon')$-evasive $k$-subspace family on $\PP^n$. Then
\[\mathcal H'=\{W\cap\Aspace^n:W\in\mathcal H,W\not\subseteq H_\infty\}\]
is an $(n,d,\epsilon)$-evasive affine $k$-subspace family on $\Aspace^n$, where $\epsilon=\epsilon'/(1-\epsilon')\le2\epsilon'$. Moreover,
\[\mathcal H''=\{W\in\mathcal H:W\not\subseteq H_\infty\}=\{W^{\rm cl}:W\in\mathcal H'\}\]
is an $(n,d,\epsilon)$-evasive $k$-subspace family on $\PP^n$.
\end{result}
\begin{proof}
By $(n,d,\epsilon')$-evasiveness of $\mathcal H$, at most $\epsilon'$-fraction of $W\in\mathcal H$ are fully contained in $H_\infty$. Throwing away those $k$-subspaces fully contained in $H_\infty$ increases the error parameter $\epsilon'$ by at most a factor of $1/(1-\epsilon')$. Therefore, $\mathcal H''=\{W\in\mathcal H:W\not\subseteq H_\infty\}$ is $(n,d,\epsilon)$-evasive. We want to prove that $\mathcal H'=\{W\cap\Aspace^n:W\in\mathcal H''\}$ is also $(n,d,\epsilon)$-evasive.
Consider a subvariety $\mathcal V\subseteq\Aspace^n$ of degree at most $d$. Let $\mathcal V_1,\ldots,\mathcal V_s$ be the irreducible components of $\mathcal V$. The projective closure $\mathcal V^{\rm cl}$ of $\mathcal V$ has the irreducible components $(\mathcal V_1)^{\rm cl},\ldots,(\mathcal V_s)^{\rm cl}$. Consider a $k$-subspace $W\in\mathcal H''$ that evades $\mathcal V^{\rm cl}$. We just need to prove that $W\cap\Aspace^n$ evades $\mathcal V$. This is true since for each $i\in[s]$,
\begin{align*}
\dim((W\cap\Aspace^n)\cap\mathcal V_i)&\le\dim(W\cap(\mathcal V_i)^{\rm cl})\\
&\le\dim(W)+\dim((\mathcal V_i)^{\rm cl})-n\\
&=\dim(W\cap\Aspace^n)+\dim(\mathcal V_i)-n
\end{align*}
where the second inequality holds since $W$ evades $\mathcal V^{\rm cl}$ and the last equality uses the fact $W\not\subseteq H_\infty$.
\end{proof}
The affine case of Theorem 1.7 now follows easily.
\begin{proof}[Proof of Theorem 1.7 in the affine case]
If $k=n$, just choose $\mathcal H=\Aspace^n$. Now assume $k<n$. Construct an $(n,d,\epsilon/2)$-evasive $k$-subspace family $\mathcal H$ on $\PP^n$ using Theorem 3.14. Then
\[\mathcal H':=\{W\cap\Aspace^n:W\in\mathcal H,W\not\subseteq H_\infty\}\]
is an $(n,d,\epsilon)$-evasive affine $k$-subspace family on $\Aspace^n$ by Lemma 3.15. The nonhomogeneous linear equations defining $W\cap\Aspace^n\in\mathcal H'$ can be easily computed from the homogeneous linear equations defining $W\in\mathcal H$ by letting $X_{n+1}=1$.
\end{proof}
The proof of Theorem 1.7 is now complete.
